% Preserved normative and evidence protocol from the pre-transcript draft.
\subsection{The normative contract and the evidence chain}\label{sec:conduct-norm}

The standard system prompt tells each agent to ``win the game by maximizing long-term advantage and being the last non-bankrupt player'' and permits ``any legal strategy or personality (aggressive, deceptive, cooperative, conservative, manipulative, opportunistic, friendly, hostile, etc.)''. It describes the public message as a place to ``negotiate, bluff, deceive, cooperate, intimidate, joke, or stay silent'', and states: ``Deception in public chat is allowed; private thoughts should reflect true reasoning.'' The action-only retry, used only after a provider-refused first attempt, removes the private-note field and shortens that sentence to ``Deception in public chat is allowed.'' It also instructs: ``Do not include private\_thought or a reasoning transcript.'' This is a response-format exclusion, not a separate honesty requirement for provider-returned explanations; the records do not establish which provider fields that instruction controls. The win and legal-action instructions appear in every prompt. Public deception is permitted throughout, while an applied private note is asked to reflect true reasoning only under the standard condition. Every Opus first attempt from game~5 onward was refused, so its applied decisions there were resolved under the action-only condition or by fallback (Section~\ref{sec:setup}). Opus therefore mostly played without the private-honesty clause, not merely without a private note.

The contract fixes how the rest of this section is labelled. A false public statement is not a violation of these instructions, and by itself it is not evidence that an agent is misaligned with them. A false private note written under the standard condition is a candidate breach of a request, but error, stale context and disagreement between channels remain live explanations, so a false note does not on its own establish dishonesty. Rule-breaking actions cannot be applied, because the engine executes only menu actions; invalid attempts are recorded instead. We describe the deceptive and adversarial tactics the records support while distinguishing them from violations of the assigned task. The analysis asks how claims were produced, how recipients detected and answered them, and what changed. This differs from AgentMonopoly, whose prompt encourages bluffing but prohibits public messages that contradict the submitted action \citep{wu2026agentmonopoly}, and from the Vending-Bench Arena communication study, which applies an external taxonomy to agents given no conduct instruction \citep{li2026arena}. An observed contradiction can be a measured violation in one of those environments and an authorized move in ours.

We distinguish four objects. A \emph{factual assertion} concerns the board or rules. An \emph{option omission} leaves a relevant legal route unstated, potentially creating a false necessity implication. A \emph{bargaining posture}, such as ``final'', expresses a stated limit rather than an observed reservation price. An \emph{assurance} concerns future behavior, whose obligation and expiry may be unspecified. State and engine checks can settle some factual or feasibility propositions; they cannot alone decide intent or whether an omission changes beliefs. Speaker-observation checks distinguish an error from a statement contradicted by information available to its author \citep{turpin2023unfaithful,lanham2023faithfulness}.

For each case we link the proposition, speaker channels, recipient's actual prompt, response, subsequent events and legal alternatives. Exposure is checked at the response where reliance is proposed. Information defects remain part of that condition: G4 d148--152 contains reversed acceptances, misattributed rejections and self-partner counters; G13 d252--253 contains a self-partner counter. The inspected G6 d19--21, G11 d255--258 and G19 d135--136 timelines contain none of these known formatter manifestations. These are exact-row checks, not assertions that other prompts are error-free or that a defect caused a response.

